% ======================================================
%  Triki zaawansowane -- hand-written.
%
%  Here the HINT is the content: the answer takes ten seconds to verify and the
%  shortcut is the thing worth owning. That is the same rule the basic volume's
%  trick set states, one level of difficulty further on.
% ======================================================

\begin{zestaw}{Triki zaawansowane}{3}{12:00}{2}
  \zz{$997 \times 998$}{995\,006 \hfill {\itshape\footnotesize Obie blisko 1000: $997-2$ \,|\, $3 \times 2$ na trzech miejscach.}}
  \zz{$87 \times 93$}{8\,091 \hfill {\itshape\footnotesize $90^2 - 3^2$.}}
  \zz{$105^2$}{11\,025 \hfill {\itshape\footnotesize $10 \times 11$, dopisz $25$.}}
  \zz{$76 \times 74$}{5\,624 \hfill {\itshape\footnotesize Te same dziesiątki, jedności sumują się do 10: $7 \times 8$ \,|\, $6 \times 4$.}}
  \zz{$1\,001 \times 437$}{437\,437 \hfill {\itshape\footnotesize $1001 = 7 \times 11 \times 13$ — powtarza trzycyfrową liczbę.}}
  \zz{$999 \times 999$}{998\,001 \hfill {\itshape\footnotesize $(1000-1)^2 = 10^6 - 2000 + 1$.}}
  \zz{$25 \times 32$}{800 \hfill {\itshape\footnotesize $\times 100$, potem $:4$.}}
  \zz{$64 \times 125$}{8\,000 \hfill {\itshape\footnotesize $125 = 1000:8$.}}
  \btnc
  \zz{Pierwiastek z $5\,329$}{73 \hfill {\itshape\footnotesize Kończy się na 9, więc pierwiastek na 3 lub 7; leży między 70 a 80.}}
  \zz{$19 \times 21$}{399 \hfill {\itshape\footnotesize $20^2 - 1$.}}
  \zz{$111 \times 111$}{12\,321 \hfill {\itshape\footnotesize Palindrom — rośnie do liczby jedynek i wraca.}}
  \zz{$1/7$ dziesiętnie, sześć miejsc}{0,142857 \hfill {\itshape\footnotesize Ten sam sześciocyfrowy cykl obraca się dla $2/7$, $3/7$ i reszty.}}
  \zz{$7/8$ na procent}{87,5\% \hfill {\itshape\footnotesize $1/8 = 12{,}5\%$, więc $7/8$ to $100 - 12{,}5$.}}
  \zz{$2^{16}$}{65\,536 \hfill {\itshape\footnotesize $2^{10} \times 2^{6}$, czyli $1024 \times 64$.}}
  \zz{Czy $4\,321 \times 1\,234 = 5\,332\,124$? \;/\; Sprawdź przez dziewiątki.}{Nie \hfill {\itshape\footnotesize Cyfry kontrolne: $1 \times 1 = 1$, a wynik daje 2.}}
  \zz{$48 \times 15$}{720 \hfill {\itshape\footnotesize $\times 10$ plus połowa tego: $480 + 240$.}}
\end{zestaw}
